% !TeX root = ../main.tex
%
% Five columns, four bands, all placement relative (positioning library). The category
% lane occupies band 1 and the region lane band 2; the region task's second tower gets
% its own band 3 so the two towers cannot touch, and the fusion node sits at the
% midpoint of the two towers in column 5 so each tower reaches it on its own line.
%
% Three label rules, one per collision the previous layout produced:
%   * the exchange label lies in the 1.5cm band gap between the two shared-context
%     nodes, so it cannot reach either node's text;
%   * the bypass label lies ABOVE its own dashed segment, not on it;
%   * every edge label carries fill=white, so an edge that passes behind it is cut.
%
% Measured width 441.59pt against this document's 455pt text block.
\begin{tikzpicture}[
    font=\small,
    node distance=0.9cm and 0.6cm,
    block/.style={draw=black!70, rounded corners=2pt, fill=black!3,
        align=center, text width=2.3cm, minimum height=1.3cm, inner sep=4pt},
    interaction/.style={block, fill=black!10, line width=0.8pt},
    head/.style={block, fill=black!7},
    output/.style={draw=black!80, rounded corners=2pt, fill=white,
        align=center, text width=2.3cm, minimum height=1.3cm, inner sep=4pt,
        double, double distance=1pt},
    flow/.style={-{Latex[length=2.2mm]}, line width=0.75pt, draw=black!80},
    privateflow/.style={-{Latex[length=2.2mm]}, line width=0.75pt, dashed,
        draw=black!70},
    exchange/.style={{Latex[length=2.0mm]}-{Latex[length=2.0mm]},
        line width=0.75pt, draw=black!75},
    edgelabel/.style={fill=white, inner sep=1.5pt, font=\scriptsize},
    group/.style={draw=black!40, rounded corners=4pt, dashed, fill=black!2,
        inner sep=6pt}
]
    % Band 1: the category lane, left to right.
    \node[block] (cin) {\textbf{Category input}\\Check-in history\\$9\times64$};
    \node[block, right=of cin] (cenc) {\textbf{Private encoder}\\$64\rightarrow256$};
    \node[interaction, right=of cenc] (cstate)
        {\textbf{Shared context}\\Category stream\\$9\times256$};
    \node[head, right=of cstate] (chead)
        {\textbf{Category head}\\Four-layer GRU\\last valid state};
    \node[output, right=of chead] (cout)
        {\textbf{Prediction}\\Next category\\seven logits};

    % Band 2: the region lane. The band gap is set wide enough for the exchange label.
    \node[block, below=1.5cm of cin] (rin)
        {\textbf{Region input}\\Region history\\$9\times64$};
    \node[block, right=of rin] (renc) {\textbf{Private encoder}\\$64\rightarrow256$};
    \node[interaction, right=of renc] (rstate)
        {\textbf{Shared context}\\Region stream\\$9\times256$};
    \node[head, right=of rstate] (shared)
        {\textbf{Shared route}\\Region tower\\shared $9\times256$ stream};

    % Band 3: the second region tower, on its own band so the towers stay apart.
    \node[head, below=of shared] (private)
        {\textbf{Private route}\\Region tower\\raw $9\times64$ history};

    % Column 5 for the region lane: fusion between the two towers, prediction below it.
    \coordinate (towermid) at ($(shared)!0.5!(private)$);
    \node[head] (fusion) at (cout |- towermid)
        {\textbf{Fusion}\\Private plus\\scaled shared};
    \node[output, below=of fusion] (rout)
        {\textbf{Prediction}\\Next region\\dataset-specific logits};

    \draw[flow] (cin) -- (cenc);
    \draw[flow] (cenc) -- (cstate);
    \draw[flow] (cstate) -- (chead);
    \draw[flow] (chead) -- (cout);

    \draw[flow] (rin) -- (renc);
    \draw[flow] (renc) -- (rstate);
    \draw[flow] (rstate) -- (shared);
    \draw[flow] (shared) -- (fusion);
    \draw[flow] (private) -- (fusion);
    \draw[flow] (fusion) -- (rout);

    \draw[exchange] (cstate) -- node[edgelabel]{two blocks} (rstate);

    % The bypass drops below the region lane and runs in the empty band-3 corridor.
    \coordinate (bypasscorner) at (rin.south |- private.west);
    \draw[privateflow] (rin.south) -- (bypasscorner)
        -- node[edgelabel, above]{raw-region bypass} (private.west);

    \begin{scope}[on background layer]
        \node[group, fit=(cstate)(rstate)] (crossgroup) {};
    \end{scope}
\end{tikzpicture}
